\section{Introdução}

O monitoramento contínuo de focos de incêndio e eventos climáticos extremos exige uma infraestrutura de dados capaz de processar e analisar grandes volumes de informações geoespaciais e temporais em tempo real. O projeto CCO-DataOps foi idealizado para atender a essa demanda, integrando dados provenientes de satélites (como os focos de calor registrados pelo INPE) com dimensões estratégicas nacionais, incluindo infraestruturas de saúde (CNES), de logística e energia (ONS), e malhas territoriais complexas (Terras Indígenas e Municípios).

Para garantir a fluidez da exploração visual dessas informações em painéis interativos (através do Apache Superset), optou-se pela utilização do Apache Druid como motor principal de banco de dados analítico (OLAP). Diferente de bancos de dados relacionais tradicionais, o Apache Druid é otimizado para ingestão particionada e compressão colunar profunda, oferecendo baixa latência na resolução de consultas de agregações matemáticas e séries temporais aplicadas a Big Data.

Entretanto, a introdução de operações espaciais e cruzamentos dimensionais (\textit{Joins}) em tempo real traz desafios computacionais severos. O tráfego de matrizes de vetores GeoJSON, atrelado à necessidade de agrupamentos multidimensionais dinâmicos (os chamados \textit{Virtual Datasets}), pode saturar rapidamente os nós de processamento da arquitetura e degradar a experiência do usuário, extrapolando a janela aceitável de latência.

Diante deste cenário, este relatório técnico apresenta um estudo analítico e exaustivo de \textit{performance} e \textit{Capacity Planning} da infraestrutura do CCO-DataOps. Através de um simulador de concorrência extrema --- que submeteu o banco de dados a diversas rodadas de estresse, perfis de hardware variados e cenários de acesso concorrente --- objetiva-se mapear o teto de escalabilidade vertical do sistema, quantificar a eficiência de respostas do \textit{Cache} de memória, decifrar o \textit{footprint} de armazenamento em disco e, primariamente, justificar cientificamente todas as decisões de projeto adotadas para sustentar a plataforma.
